\newpage
\fancyhead[L]{\songti\zihao{-5}\cquHeader}
\fancyhead[R]{\songti\zihao{-5}致谢}
\addcontentsline{toc}{section}{致谢}
\section*{致\quad 谢}
\zihao{-4}

潺潺流水，终会穿过一座座群山。

感谢这一路途中遇见的所有，它们不断地推着我往前走。

感谢遇见的所有老师。

桃李不言，下自成蹊。教师是一份非常纯粹而辛苦的职业，需要不断向外输出、不断给予，且常常不求回报。我很幸运，在重庆大学遇见的每一位老师，都最大程度地帮助着包括我在内的许多学生，这份纯粹令我敬佩。

特别感谢赵宏博老师。在大一、大二的时候，我从赵老师的课程和知乎回答中受益良多，因此主动选择了赵老师的论文课题。同时，也感谢他在本论文以及过往课程中给予的所有指导。赵老师可以看作学院内许多老师的缩影：善良且纯粹，认真却不严苛。这种态度真真切切地影响了我人格的塑造。也许此生无缘再见，但过往的烟云也无法堙没如金子般的教诲。

感谢重庆以及重庆大学。

四年里，我在学校的庇护下，得以较为从容地学习、生活、犹疑，也得以在尚未成熟时保有一点自由。师如道公，我如顽石。我未能成为足够优秀的学生；面对学校无条件的包容，而我所能回报的实在太少。念及此处，惭愧之外，唯有感恩。

“自己选择，自己承担”，这是在学校里我最喜欢的一句话。它朴素，却有分量；平静，却依然坚毅。它也在我面对人生无数选择时给予我力量。

感谢数学。

高山仰止，景行行止。很遗憾，我没能为这个学科做出多少贡献，但余虽愚，亦有所闻。如今再回头看，在学习过程中，学过哪些内容、记住多少定理，反倒没有那么重要。真正留下来的，是所谓的 \emph{mathematical maturity}：一种愿意静心思考的习惯，一种面对复杂问题时不轻易放弃的耐心，一种保持纯粹、保持分析、保持追问的能力。这种思维方式，是数学给予我最宝贵的东西。

若从现实角度来看，数学似乎并不是一门“实用”的专业。它很少有直接对口的岗位，也常常需要再补充许多其他学科的知识，才能进入具体的行业之中。但以结果来看，学数学的人大多也都走得不错。究其原因，也许正是因为数学慷慨地把那种愿意动脑、愿意想办法、愿意在混沌中寻找秩序的能力，赠予了每一个曾认真学习它的人。对我而言，这已经足够珍贵。

我最终未必会在以后继续学习数学，但仍感谢它曾以沉默而严格的方式塑造过我。此后无论去往何处，这份训练与馈赠，都将长久地留在我身上。

从小县城走到大学，又从大学走到社会，我所面对的世界越来越大，不变的是我从中感受到的善意。

感谢这个时代与社会所提供的机会，让我能探索自己的人生。感谢繁荣昌盛的祖国能够供养我完成学业，使一个从小县城走出来的普通学生，能够通过读书看见更大的世界。感谢学校和社会那些陌生但帮助过我的人。感谢身边一直支持我、给予我温暖的人。

我会时刻铭记并传递这同一份善意。

日子无论怎样慢慢流逝，终究也过去了。以自我成长来说，这四年无疑是人生最深刻的四年：也彷徨，也失路，辗转反侧，如梦初醒。也感谢那个如潺潺流水般，终究穿过一座座群山的自己。

最后，感谢为本论文审查、评阅和答辩付出辛勤劳作的各位老师。
